\documentclass{autre}

\title{Analogies entre les différentes grandeurs transportée}
\date{}

\begin{document}
  \maketitle{}\thispagestyle{empty}
  \setlength{\extrarowheight}{40pt}
  \renewcommand{\tabularxcolumn}[1]{>{\small}m{#1}}
  \newcolumntype{C}{>{\centering\arraybackslash}X}
  \noindent
  \begin{tabularx}{\textwidth}{| X | C | C | C | C | C | C |}
    \hline
    Grandeur conservative &
    Charge &
    Chaleur &
    Particules &
    Masse &
    Volume &
    Quantité de mouvement suivant $\vv{u}_x$\\
    \hline
    Unité de la grandeur & 
    &
     &
     &
     &
     &
    \\
    \hline
    \vspace{1em}Grandeur par unité de volume et son unité\vspace{1em} &
     &
     &
     &
     &
    \cellcolor{gray} &
    \\
    \hline
    Vecteur densité de courant et son unité &
     &
     &
     &
     &
     &
     \\
    \hline
    Débit et son unité &
     &
     &
     &
     &
     &
     \\
    \hline
    Équation locale de conservation de la grandeur et son origine &
     &
     &
     &
     &
    \cellcolor{gray} &
     \\
    \hline
    Équation de conservation \textbf{en régime stationnaire} (locale et globale) &
    &
    &
    &
    &
    &
    \cellcolor{gray} \\
    \hline
    Vecteur densité de courant comme grandient &
     &
     &
     &
    \cellcolor{gray} &
    \cellcolor{gray} &
     \\
    \hline
    Équation de diffusion &\cellcolor{gray} &
     &
     &
    \cellcolor{gray} &
    \cellcolor{gray} &
    \cellcolor{gray} \\
    \hline
  \end{tabularx}
\end{document}
